% !TeX root = ../main.tex

\section{Имитационное моделирование}
\label{sec:simulation}

\subsection*{Общая схема генерации и анализа данных}

Полные нелинейные ответы P-DCM и BW рассчитывались при известных амплитудах входа. Для каждого заданного расписания генерировался BOLD-сигнал, по нему оценивались коэффициенты ОЛМ с канонической HRF, после чего вычислялись одношаговые поправки. Оценки сравнивались с заданными амплитудами $a_A$, $a_B$ и их разностью $\Delta$. В бесшумных контрольных расчетах дополнительно использовалось собственное линейное ядро $L_1$. Параметры генераторов, канонической HRF и поправок не подбирались по ошибкам тестовых оценок.

Каждое событие представляло собой прямоугольный импульс длительностью 1~с. Интервал между началами последовательных событий (stimulus onset asynchrony, SOA) задавался как
\[
t_{i+1}-t_i=1+\mathrm{Exp}(\overline{\mathrm{SOA}}-1),
\]
где $\mathrm{Exp}(b)$ означает экспоненциальную случайную величину со средним $b$ секунд. Такая схема исключает перекрытие самих стимулов, сохраняя случайное перекрытие гемодинамических ответов. Моменты начала округлялись до ближайшего узла сетки $0{,}05$~с; первое событие начиналось на пятой секунде записи. Основные средние SOA составляли 2, 4 и 8~с; при испытании статической поправки дополнительно использовался SOA 16~с. Окно стимуляции составляло 480~с, за ним следовали 64~с покоя, так что полная длительность ряда равнялась 544~с. Период регистрации (repetition time, TR) составлял 1~с. Времена предъявления не оптимизировались под эффективность оценок~\cite{wager2003}.

В основном режиме каждому событию независимо назначалось условие A или B с вероятностью $0{,}5$. Для исследования влияния истории использовались строгое чередование и метки, связанные с предшествующим интервалом: событие относилось к A, если интервал от предыдущего события не превышал медиану интервалов данного расписания, и к B в противном случае. При сравнении методов расписание, амплитуды и реализация шума были общими. При сравнении генераторов сохранялись также физический вход и общий E/I-блок.

Уравнения интегрировались методом Рунге--Кутты четвертого порядка с шагом $0{,}005$~с; сигнал сохранялся каждые $0{,}05$~с. Уменьшение шага до $0{,}0025$ и $0{,}00125$~с меняло ошибку контраста на величину порядка $10^{-13}$. Квадратичный оператор для полной поправки вычислялся на сетке лагов с шагом $0{,}05$~с в окне 64~с. Все расчеты начинались из состояния покоя; форма и масштаб канонической HRF были общими для серий. Дополнительная нормировка каждого BOLD-ряда на его максимум не выполнялась.

В шумовых испытаниях к сигналу добавлялся белый шум либо его смесь с процессом Орнштейна--Уленбека с постоянной времени 2~с и долей в дисперсии $0{,}5$. Отношение сигнала к шуму (signal-to-noise ratio, SNR) определялось как отношение пика BOLD-ответа номинальной P-DCM на одиночное секундное событие при $a=0{,}2$ к стандартному отклонению шума и равнялось 2 или 8. Для оценивания использовалась известная ковариация соответствующего шума.

Серии решали разные задачи и анализировались раздельно (табл.~\ref{tab:series}). Сводные оценки полной и статической поправок относятся к разным наборам условий и не служат их прямым сравнением на единой выборке.

\subsection*{Амплитуда и плотность событий}

Первая серия проверяла, какая ошибка остается при конечных амплитудах и как она зависит от плотности событий. Использовались девять случайных последовательностей, по три для каждого среднего SOA 2, 4 и 8~с. Для каждой последовательности задавались пять значений $a_A$: $0{,}025$, $0{,}05$, $0{,}1$, $0{,}2$ и $0{,}3$; амплитуда второго условия составляла $a_B=a_A/2$. Для обоих генераторов рассчитывались бесшумные ответы, а амплитуды и контраст оценивались с $C$ и собственным $L_1$. Таким образом, сравнивались 90 сочетаний генератора, расписания и пары амплитуд.

\subsection*{Общий уровень стимуляции и история условий}

Во второй серии использовалась та же база из девяти расписаний. В первой постановке разность амплитуд удерживалась равной $\Delta=0{,}1$, а обе амплитуды увеличивались: $(0{,}15;\,0{,}05)$, $(0{,}20;\,0{,}10)$, $(0{,}25;\,0{,}15)$ и $(0{,}30;\,0{,}20)$. Это позволяло менять средний уровень $m=(a_A+a_B)/2$ при постоянном контрасте.

Во второй постановке амплитуды были равны: $(0{,}1;\,0{,}1)$, $(0{,}2;\,0{,}2)$ и $(0{,}3;\,0{,}3)$. При таких входах полный сигнал не зависит от назначения меток A и B. Для каждого расписания сравнивались случайные метки, строгое чередование и метки по истории интервалов; дополнительно выполнялись 100 контрольных перестановок. Ненулевая оценка при $\Delta=0$ характеризовала влияние разметки и предыстории на проекцию нелинейного остатка. Все ответы рассчитывались для обоих генераторов без шума.

Отдельно исследовались одиночные события и пары одинаковых событий при $a=0{,}05$, $0{,}15$ и $0{,}30$ и SOA 2, 4, 8 и 16~с. Полный расчет сопоставлялся со вторым порядком по интегралу взаимодействия пары и по остатку одиночного ответа, деленному на $a^2$. Эти контролируемые воздействия позволяли разделить отклонение отдельного ответа от линейного масштабирования и взаимодействие соседних событий.

\subsection*{Испытание полной поправки Вольтерра}

Третья серия проверяла чувствительность фиксированной поправки FV к параметрам и структуре генератора. Рассматривались три режима: номинальная P-DCM, BW и 40 вариантов P-DCM с параметрами, выбранными латинским гиперкубом. В последнем режиме время среднего транзита крови $\tau_0$ менялось от 1 до 5~с, постоянная вязкоупругого оттока $\tau_v$ от 0 до 16~с, а параметр наблюдения $\epsilon_{\mathrm{obs}}$ от $0{,}1291$ до $0{,}5648$, изменяя коэффициенты $c_2$ и $c_3$ в~\eqref{eq:a3}. Во всех трех режимах использовалась одна поправка, построенная по номинальной P-DCM.

Испытание включало 120 независимых случайных последовательностей, по 40 на каждый средний SOA 2, 4 и 8~с. Задавались восемь пар амплитуд: $(0{,}05;\,0{,}025)$, $(0{,}1;\,0{,}05)$, $(0{,}2;\,0{,}1)$, $(0{,}3;\,0{,}15)$, $(0{,}3;\,0{,}2)$, $(0{,}1;\,0{,}1)$, $(0{,}3;\,0{,}3)$ и $(0{,}2;\,0{,}18)$. Первые пять пар образовывали основной класс ненулевых контрастов, последняя пара задавала малый контраст $\Delta=0{,}02$, две пары с равными амплитудами задавали нулевой контраст. Использовались бесшумные ответы и два уровня шума, SNR 8 и 2, с восемью реализациями шума.

Форма поправки была выбрана на отдельной проверочной выборке до независимого испытания и после его открытия не менялась. Для представленных трех режимов сводная RMSE рассчитывалась с равными весами 27 групп: трех режимов генератора, трех классов контраста и трех средних SOA. Внутри групп сохранялись равные веса пар амплитуд и шумовых условий. Изменение ошибки определялось как $H_R=R_{\mathrm{FV}}-R_C$. Для оценки применимости использовались заданные до испытаний допуски: снижение сводной RMSE не менее $0{,}002$ и увеличение RMSE отдельной группы не более $0{,}001$ модельной единицы. Сужение представляемой области не сопровождалось новым выбором поправки.

\subsection*{Испытание статической поправки винеровского типа}

Четвертая серия проверяла, при каких условиях FW снижает ошибку до заданного уровня. Для каждого генератора, P-DCM и BW, использовались собственная поправка, рассчитанная по той же модели, и перенесенная поправка, рассчитанная по другой модели. Эти четыре сочетания сравнивались при средних SOA 2, 4, 8 и 16~с. Для каждого SOA и режима назначения меток генерировалось 47 случайных последовательностей; использовались случайные метки и метки по истории интервалов.

Задавались 20 пар амплитуд: для каждого $a_A\in\{0{,}025;\,0{,}05;\,0{,}1;\,0{,}2;\,0{,}3\}$ использовались отношения $a_B/a_A\in\{0;\,0{,}5;\,0{,}9;\,1\}$. Амплитуды объединялись в три уровня: малые ($a_A=0{,}025$ и $0{,}05$), средние ($0{,}1$ и $0{,}2$) и большие ($0{,}3$). Бесшумные ответы и уровни SNR 8 и 2 образовывали три шумовых условия; на каждом ненулевом уровне SNR учитывались белый и коррелированный варианты шума с восемью реализациями. Сочетания четырех вариантов генератора и поправки, четырех SOA, трех уровней амплитуды и трех шумовых условий образовывали 144 группы. Испытание выполнялось с фиксированными настройками.

Для каждой группы задавался допуск $\delta_s=\max(0{,}002;\,0{,}05m_s)$, где $m_s$ равен среднему модулю заданного контраста либо характерной амплитуде при оценке коэффициентов. Исход <<исходная оценка достаточна>> фиксировался, если верхняя граница ошибки $C$ не превышала $\delta_s$. Исход <<снижение ошибки подтверждено>> требовал уменьшения ошибки более чем на $0{,}002$ при увеличении не более чем на $0{,}001$ модельной единицы. Улучшение считалось полным, если после поправки ошибка не превышала $\delta_s$, и частичным в противном случае. Если ни $C$, ни FW не достигали заданной точности, фиксировался соответствующий исход; при недостаточной точности интервалов исходы считались неразличимыми.

\subsection*{Пределы второго порядка и подстановка оцененных амплитуд}

В пятой серии сумма линейного и квадратичного ответов сравнивалась с полным расчетом обоих генераторов при SOA 2, 4 и 8~с. Основная последовательность $k=0$ использовалась для пяти амплитуд при $a_B=a_A/2$, двух пар с постоянным контрастом $(0{,}1;\,0)$ и $(0{,}3;\,0{,}2)$, а также равных амплитуд $0{,}1$ и $0{,}3$ при случайной и связанной с историей разметке. Дополнительно проверялась пара $(0{,}3;\,0{,}15)$ с метками по истории. Последовательность $k=1$ служила контролем для пары $(0{,}3;\,0{,}15)$ и равных амплитуд $0{,}3$. Всего сравнивались 90 сочетаний. FV и FW вычислялись с заданными амплитудами и с подстановкой $\hat\beta_C$; первый вариант является теоретическим контролем, недоступным в реальных данных. Эти бесшумные расчеты не объединялись с независимыми испытаниями поправок.

Для оценки согласия второго порядка с полным расчетом знак ошибки считался различимым, если ее модуль в обоих расчетах был не менее $0{,}002$. Величины считались близкими при расхождении не более $25\%$. Различия ошибок менее $2\cdot10^{-5}$ модельной единицы считались численно незначимыми.

\subsection*{Иллюстративные записи}

Для наглядного сопоставления полного BOLD-ответа, линейной модели и поправок использовались две отдельные бесшумные записи P-DCM длительностью 70~с. В первой задавались $a_A=0{,}30$, $a_B=0{,}05$ и средний SOA 1,7~с со случайными метками. Во второй при $a_A=a_B=0{,}30$ использовались средний SOA 1,8~с и метки по истории интервалов. Моменты событий генерировались в окне до 62~с, начиная со второй секунды; сохранялись стимулы, завершавшиеся до 66~с. Это отдельные примеры искажения оценок, они не включались в сводные метрики и не использовались для оценки частоты ошибок.

\subsection*{Неопределенность и контроль вычислений}

Приближенные интервалы получены бутстрепом по последовательностям стимулов (1000 выборок) с нормальной аппроксимацией и поправкой Бонферрони при $\alpha=0{,}05$. Для выделенных режимов FV сохранены исходные верхние границы: семейство сравнений после сужения области не уменьшалось. Ключевые вычисления, включая производные и ядра, контроль интегрирования, случайные выборки и основные агрегаты, ранее продублированы независимой реализацией. Эта проверка не заменяет оценку физиологической применимости генераторов.

\begin{table}[!ht]
\caption{Имитационные серии и соответствующие результаты}
\label{tab:series}
\centering
\small
\setlength{\tabcolsep}{4pt}
\begin{tabular}{@{}p{0.065\linewidth}p{0.425\linewidth}p{0.29\linewidth}p{0.15\linewidth}@{}}
\hline
Серия & Вопрос & Данные & Результат\\
\hline
1 & Амплитуда и плотность событий & 9 расписаний, 5 пар, два генератора; без шума & Рис.~\ref{fig:f5}, табл.~\ref{tab:amp}\\[3pt]
2 & Общий уровень, история и взаимодействия событий & Те же 9 расписаний; одиночные и парные события; без шума & Рис.~\ref{fig:f3}, \ref{fig:f6}\\[3pt]
3 & Полная поправка при несоответствии модели & 120 расписаний, 8 пар, три режима; без шума, SNR 8 и 2 & Рис.~\ref{fig:f7}\\[3pt]
4 & Достижение заданной точности статической поправкой & 47 расписаний на SOA и режим меток, 20 пар, 144 группы & Табл.~\ref{tab:fw}\\[3pt]
5 & Второй порядок и подстановка оценок & Оба генератора, заданные и оцененные амплитуды; без шума & Рис.~\ref{fig:f8}\\
\hline
\end{tabular}
\end{table}
